\section{Discussione}\label{sec:discussion}

Le gravi discrepanze evidenziate dai test statistici sottolineano i limiti dell'affidarsi esclusivamente a registri amministrativi volontari per quantificare i fenomeni demografici.
\subsection{Fattori Umani e Amministrativi dietro la Sottostima}
Il bias sistematico osservato nel test di Wilcoxon — una costante mancanza di registrazioni AIRE rispetto ai dati dei paesi di destinazione — è guidato da un disallineamento degli incentivi.
I cittadini che emigrano trovano ragioni pratiche immediate per registrarsi nel paese di destinazione (per ottenere casa, lavoro, assistenza sanitaria e credenziali fiscali), garantendo un conteggio molto accurato da parte delle istituzioni straniere (Eurostat).
Al contrario, registrarsi all'AIRE italiana fornisce scarsi benefici immediati al cittadino e richiede un impegno proattivo con le procedure burocratiche.
Di conseguenza, molti emigranti ritardano la registrazione per anni o la omettono del tutto.
Questo ritardo spiega anche i risultati del test del Chi Quadrato, che ha dimostrato come la distribuzione temporale dei dati ISTAT non corrisponda ai flussi migratori reali.
Un individuo che emigra fisicamente in Germania nel 2017 potrebbe registrarsi all'AIRE solo nel 2019, quando necessita di un servizio consolare (come il rinnovo del passaporto).
Questo ritardo di fase scollega l'evento amministrativo dalla realtà demografica, distorcendo le tendenze temporali nei dati.
\subsection{L'anomalia francese del 2020-2021}
L'unica eccezione alla sottostima sistematica si è verificata con la Francia durante la pandemia di COVID-19 (2020–2021), dove l'ISTAT ha registrato più emigranti di Eurostat.
Questo "gap negativo" è probabilmente il risultato di due fattori che si intersecano.
In primo luogo, la vicinanza e la natura altamente circolare della migrazione IT-FR hanno fatto sì che molti italiani residenti in Francia senza registrazione formale potrebbero essere tornati bruscamente in Italia durante i primi lockdown, o viceversa, potrebbero aver formalizzato il proprio status a causa delle incertezze legate alla pandemia.
In secondo luogo, a differenza dei rigidi sistemi di registrazione obbligatoria di paesi come la Germania (\textit{Anmeldung}), il sistema francese per la registrazione dei cittadini UE può essere meno rigido, portando potenzialmente a un conteggio inferiore o ritardato sul lato francese durante le interruzioni amministrative della pandemia.
\subsection{Implicazioni Politiche}
L'incapacità del registro AIRE di tracciare accuratamente la scala e la tempistica dell'emigrazione italiana ha profonde implicazioni.
Le decisioni politiche, le previsioni sul mercato del lavoro e le strategie di investimento regionale si basano attualmente su dati che sottostimano significativamente la fuga di capitale umano.
I nostri risultati suggeriscono che ricercatori e decisori politici devono integrare sistematicamente le statistiche speculari delle nazioni riceventi per correggere i punti ciechi della rendicontazione nazionale.
